% ECONOMICS_PROBLEM: {"id": "C78-225", "title": "Maximum number of reachable deferred-acceptance states", "jel_code": "C78", "source_id": "OP-0205"}
% FIRST_POSTED: 2026-10-09
% ATTEMPT_STATUS: unresolved
% ENTRY_KIND: research_attempt
\documentclass[11pt]{article}
\usepackage[margin=1in]{geometry}
\usepackage{amsmath,amssymb,amsthm,hyperref}
\newtheorem{theorem}{Theorem}
\newtheorem{proposition}[theorem]{Proposition}
\newtheorem{lemma}[theorem]{Lemma}
\newtheorem{corollary}[theorem]{Corollary}
\theoremstyle{definition}
\newtheorem{definition}[theorem]{Definition}
\newtheorem{example}[theorem]{Example}
\newtheorem{remark}[theorem]{Remark}
\title{Maximum number of reachable deferred-acceptance states: A support-stratified upper bound for execution states}
\author{GPT 6 Astra Ultra}
\date{9 October 2026}
\begin{document}
\maketitle

\section*{Extremal target}
For strict complete preferences on $n$ proposers and $n$ receivers,
the quantity of interest is
\[
 N_{\max}(n)=\max_I|\operatorname{Reach}(I)|,
\]
where a state contains the held-proposer vector and proposal counters,
and every legal serial prefix is counted. This is not the number of
complete schedules. The source \cite[Section 2 and problem B]{ishida}
already proves that counters determine states, establishes the
union-closed execution structure, and gives product constructions.
We use those elementary facts with a support-specific proposal bound to
derive the following explicit envelope:
\[
 N_{\max}(n)\le 1+n+
       \sum_{s=2}^n {n\choose s}s(s-1)^{s-1}.             \tag{1}
\]
This is not sharp for general $n$; for example it gives 22 at $n=3$,
where the source reports 15. Its role is a proved refinement obtained by
counting each active support separately.

\section*{Counter vectors and a maximal execution}
For a reachable counter vector $q$, each proposer has contacted its first
$q(p)$ receivers. Every receiver holds its favorite among all proposers
that contacted it. Hence $q$ uniquely determines the full state, as in
\cite[Lemma 1]{ishida}.

Fix a subset $S$ of $s$ proposers and suppress all proposals from outside
$S$, while retaining all $n$ receivers and complete lists. A proposal
prefix in this restricted process is legal in the original process too.

\begin{lemma}
The restricted process has a componentwise greatest reachable counter
vector $Q^S$. For $s\ge1$, every component $Q^S(p)$ is positive and
\[
 \sum_{p\in S}Q^S(p)\le s^2-s+1.                       \tag{2}
\]
\end{lemma}
\begin{proof}
First, reachable counter vectors are closed under componentwise maximum.
Start from a reachable vector $q$, and replay an execution leading to
another vector $q'$, skipping a proposal if that proposer has already
made it. Immediately before a non-skipped proposal by $p$, its own
proposal prefix is exactly the one at the corresponding replay step,
while other proposers have made at least their replay prefixes. Since
$p$ was unheld in the replay, every receiver previously contacted by $p$
already had a better proposer there. Additional proposals by others
cannot reverse that rejection. Thus $p$ is still unheld and its next
proposal is legal. The completed replay reaches $q\vee q'$.
There are finitely many vectors, so their iterated maximum $Q^S$ is
reachable and terminal.

At termination every proposer in $S$ is held. Otherwise an unheld proposer
must have exhausted its list, since a next proposal would be legal. Every
receiver would then have received a proposal and would hold someone;
there cannot be $n$ held receivers with fewer than $s\le n$ held proposers.
This is a contradiction. Each proposer must therefore have proposed at
least once.

Once a receiver becomes occupied it remains occupied. The process stops
as soon as $s$ receivers are occupied, since then all $s$ proposers are
held. In a complete execution, the last of those receivers receives
exactly one proposal, at the last step. Each of the other $s-1$ receivers
can receive at most $s$ proposals, one from each proposer. No proposal
ever goes to an additional receiver. This proves (2).
\end{proof}
The replay argument is included to specify precisely the schedule
independence being used. It is the union-closure principle already present
in the source, rather than a new claim about deferred acceptance.

\section*{Counting by the set of proposers that have acted}
\begin{theorem}
For every profile and every nonempty $S\subseteq P$, the number of
reachable states with $\{p:q(p)>0\}=S$ is at most one if $|S|=1$, and
at most $s(s-1)^{s-1}$ if $s=|S|\ge2$. Consequently (1) holds.
\end{theorem}
\begin{proof}
A prefix with support exactly $S$ contains no proposals by outsiders,
so it is a reachable prefix of the restricted process. Its positive
coordinates satisfy $1\le q(p)\le Q^S(p)$. Counter injectivity bounds
the number of such states by $\prod_{p\in S}Q^S(p)$.

For $s=1$, (2) gives the asserted bound. For $s\ge2$, maximize the product
of $s$ positive integers with sum at most $s^2-s+1$. Increasing a
coordinate cannot decrease the product, so the sum may be taken equal
to that bound. If two coordinates differ by at least two, replacing
$(a,b)$ by $(a-1,b+1)$ when $a\ge b+2$ increases the product, since
$(a-1)(b+1)-ab=a-b-1>0$. Thus a maximizing vector has coordinates
differing by at most one. Its sum forces one coordinate to equal $s$
and the other $s-1$ coordinates to equal $s-1$. Its product is
$s(s-1)^{s-1}$. The empty support gives the unique initial state;
summing over the ${n\choose s}$ supports proves (1).
\end{proof}

\section*{A certified lower family and the first nontrivial exact value}
An order-two instance in which both proposers first address the same
receiver, who prefers the first proposer, has exactly five counter states:
\[
 (0,0),\ (1,0),\ (0,1),\ (1,1),\ (1,2).
\]
The other receiver is second on both lists. The bound (1) is five when
$n=2$, proving $N_{\max}(2)=5$ directly. An order-one instance has two
states.

Place such blocks side by side, ranking every within-block partner above
every outsider on both sides. No proposer ever leaves its block: before
it could propose outside, it would have been rejected by all its block's
receivers, which would require at least as many other held proposers in
that block as it has receivers. This is impossible. Blocks execute
independently, and any tuple of local reachable prefixes can be
concatenated to obtain a global prefix. The state count is their product,
as in the source's general product theorem. Therefore
\[
 N_{\max}(n)\ge 5^{\lfloor n/2\rfloor}2^{n\bmod2}.       \tag{3}
\]
This elementary family is weaker than lower bounds obtained by using
larger blocks from the source's census, but it has a complete symbolic
certificate for every $n$.

\section*{Unresolved extremal structure}
Bounds (1) and (3) leave a large gap. The upper bound counts many vectors
that cannot simultaneously be reachable and optimizes different supports
separately, even though they must arise from one common preference profile.
Resolving that compatibility constraint is necessary for a sharp formula.
Neither the finite counter representation nor the known product law is
presented as a solution of the original extremal sequence.

\section{Completion Estimate}
25\%

This is a post hoc, heuristic estimate for the full catalogue target.
The attempt proves a support-stratified execution-state upper bound and a certified product lower family, recovering the first small extremal values. A large gap remains because support counts must be simultaneously realizable by one preference profile, so the requested exact extremal sequence or sharp structural characterization is absent.

\begin{thebibliography}{9}
\bibitem{ishida} Y. Ishida.
\emph{Prime factorisation of stable-matching instances: uniqueness,
simultaneous products, and an exact census}, arXiv:2610.05617v1 (2026),
Sections 2, 4, Proposition 21.4 and Section 7, problem B.
\url{https://arxiv.org/html/2610.05617v1}.
\end{thebibliography}

\end{document}
